% !TEX root = ../math.tex
\chapter{Stimulus Optimization: Ideal Timer Placement}
\label{chap:optimize-stimulus-parameters}

\section{The Continuous One Bit Channel}
\subsection{Two Interpretations}

The continuous one bit channel can be conceived in two ways: as a continuous channel
with a one bit range, or as a series of times of identical events. The equivalency follows
from the fact that the continuous channel may be characterized by the times at which
its signal changes. To have a finite (differential) capacity we must constrain the channel to a minimum duration of signal,
or equivalently, a minimum time between events. For our purpose, we prefer the latter interpretation of
the one bit channel, since the user communicates using a sequence of identical events such as blinks
or taps.

\subsection{The Ideal Timing Distribution}
What then, is the capacity of the one bit channel? We prefer the event interpretation for analysis.
We consider the ideal distribution of the duration between events.
Let $T$ be the minimum duration between events and $X$ be the additional delay beyond the minimum,
so that the total duration between events is $T + X$.
Let $A = T + E[X]$ be the expected inter-arrival time.
The rate is given by
$$C = \max_{\mathcal{P}_X} \frac{h(X)}{A}.$$

We refer to the rate-maximizing distribution of $X$, $\mathcal{P}_X$, as the ideal timing distribution.
We will proceed to show that the ideal timing distribution is an exponential distribution.

We treat the capacity, expected value, and differential entropy as functionals of the probability density function $f(x)$.
To find the optimal distribution, we introduce a Lagrange multiplier $\mu$ for the normalization constraint $\int f(x) dx = 1$ and define the functional $J(f)$:
$$J(f) = \frac{h(X)}{A} - \mu \left( \int f(x)dx - 1 \right)$$
We maximize $J(f)$ by setting the variational derivative with respect to $f(x)$ to zero:
$$\frac{\delta J}{\delta f(x)} = 0$$
Using the quotient rule for variational derivatives, and noting that $\frac{\delta h(X)}{\delta f(x)} = -1 - \ln f(x)$ and $\frac{\delta A}{\delta f(x)} = \frac{\delta E[X]}{\delta f(x)} = x$, we have:
\[
\frac{\delta}{\delta f(x)} \left( \frac{h(X)}{A} \right) = \frac{A(-1 - \ln f(x)) - h(X)x}{A^2}
\]
Substituting this into the stationarity condition:
\[
\frac{A(-1 - \ln f(x)) - h(X)x}{A^2} - \mu = 0
\]
Solving for $f(x)$:
\begin{align*}
f(x) &\propto \exp\left( - \frac{h(X)}{A} x \right)
\end{align*}
proving that $X$ is exponentially distributed with rate parameter equal to the rate of the channel
(i.e., $\lambda^* = C$).
This fact lets us solve for the rate,
$$C = \frac{1 - \ln C}{T + 1/C}$$
which by algebra has the solution,
$$C = W(T) / T$$, where $W$ is the Lambert W function.
% TODO: Technically the argument to W should be dimensionless
% and we should write W(T/Δ)
\subsubsection{A Geometric Intuition}
The above result (that the ideal timing distribution is exponential) may
be appreciated more readily with a geometric argument.
We consider a large sample, $\mathbf{X}$, in $N$-dimensional space,
representing the sequence of additional delays $X_i$.
By the law of large numbers, its $L_1$ norm, $\|\mathbf{X}\|_1 = \sum_i X_i$, must be
tightly distributed about the $N$-simplex, $\|\mathbf{X}\|_1 = NE[X]$.
This can be thought of as the permissible region for $\mathbf{X}$. Thus the entropy maximizing
distribution for $\mathbf{X}$ should assign a uniform distribution over the $N$-simplex,
which is precisely what the exponential distribution does as it
depends only on the $L_1$ norm, $f(\mathbf{X}) \propto \exp(-\lambda \|\mathbf{X}\|_1)$.

(Incidentally, this conception also provides a geometric proof of the fact that the
differential entropy of the exponential distribution is $1 - \ln \lambda $, since 
the entropy of a uniform distribution over the $N$-simplex should be the logarithm of
its volume. Since its side length is $N/\lambda$, its volume is $(N/\lambda)^N / N!$,
and its entropy per sample (by Stirling's approximation) is
$$h(X) = \lim_{N \to \infty} \frac{1}{N} \left( N (\log(N - \log \lambda) - N (\log N - 1)) \right) = 1 - \ln \lambda$$
as expected.)




\section{Periodic Signals}
% Discuss maximizing the rate in practice
\subsection{Random Scan Times}
\subsection{Periodic Signals}
% Discuss the advantages of making D a periodic variable
% - velocity is constant / doesn't need to be estimated (though this could be avoided)
% - info gain does not depend on the distribution of T
\subsection{Approximating a Uniform Distribution with a Gaussian Mixture}

If every prefix is assigned a periodic timer, with
a period of $P$, and a phase of $\phi_n$, then the phase
of the user's signal will be a mixture of wrapped Gaussians,
$$
P(x) = \sum_{n=1}^N \pi_n \mathcal{W}(x \mid \phi_n, \sigma^2)
$$
where $\mathcal{W}$ is the wrapped normal distribution,
and $\sigma^2$ is a constant, describing the user's
imprecision in timing their signals.
Our goal is to set the phases of the timers, $\phi_n$,
to maximize the entropy of this mixture.

We will perform this routine after the user makes a signal,
and before showing the new timers. Hence this computation
should run on the order of a millisecond or less.

Our approach is to optimize (with a quasi-Newton method)
an approximation to the entropy of the mixture
in the frequency domain.

First, we choose to use the collision entropy to
approximate the Shannon entropy of the mixture.
This is justified because, near the uniform distribution,
both objectives are maximized at the same point, and
their leading deviation from the optimum is proportional
to the same quadratic functional, namely the squared
$L^2$ distance from uniform. Thus, in the regime relevant
here, maximizing Shannon entropy is well
approximated by maximizing collision entropy.
This objective of maximizing $-\log \int P(x)^2 dx$
is equivalent to minimizing the power, $\int P(x)^2 dx$.
We will choose to minimize the power in the frequency domain.

The Fourier expansion of a wrapped normal distribution
centered at $0$ is given by
$$
a_m = \frac{1}{P} \exp\left(-\frac{2\pi^2 m^2 \sigma^2}{P^2}\right)
$$
Thus the coefficients of the mixture are,
$$
C_m = a_m s_m,
$$
where $s_m$ are the interference coefficients, given by
$$
s_m = \sum_{n=1}^N \pi_n \exp(i 2\pi m \phi_n / P).
$$
The objective function to minimize is then,
$$
J = \sum_{m=1}^F |C_m|^2
$$
where $F$ is the number of terms in the Fourier expansion.
Because the terms, $a_m$, fall off rapidly, we may
truncate to a small number of terms, $F$. For a given precision,
$\epsilon$, we may choose $F$ so that the first omitted coefficient satisfies
$$
(a_m)^2 < \epsilon.
$$
For reasonable values of $P = 1.0$, $\sigma = 0.020$,
and $\epsilon = 10^{-6}$, a value of $F = 30$ suffices.
(In the current implementation we always use the hardcoded constant $F = 32$.)
This justifies our choice to calculate the power in the frequency domain.
It is not necessarily computationally cheaper per term relative to the time domain,
but we may use far fewer terms, $F$, than the number of discretized bins
that would be required in the time domain.

In our implementation, we then optimize the phases against this objective
starting from initial equally spaced phases, using
the L-BFGS algorithm. To reduce latency, we stop before convergence,
at a fixed number of iterations, $I$ (in our implementation, $I = 25$).
This allows our implementation to run over $N=300$ components in $500\mu$s.

These choices of parameters are illustrated in figure \ref{fig:tplace-sample-plot},
which shows visible artifacts when the number of Fourier terms is too small
($F=16$), and no discernible difference between $F=32$ and $F=100$.
It also shows that there is little improvement in differential entropy after
$I=25$ iterations of L-BFGS.


\begin{figure}[H]
    \centering
    \includegraphics[width=\textwidth]{figures/tplace_sample_plot.png}
    \caption{
        \textbf{Phase Optimization for Different Values of $F$ and $I$.}
        This figure shows a wrapped normal mixture distribution with
        $N=300$ components and synthetic weights, $\pi_n$, taken
        as the softmax of a normal distribution.
        The top plot, "Equally Spaced Phases", shows the mixture if we
        assign phases as $\phi_n = n / N$.
        The lower plots show the mixture after the optimization of $\phi_n$,
        for different values of $F$ and $I$.
    }
    \label{fig:tplace-sample-plot}
\end{figure}


\subsubsection{Selection\_Mass}
The mixture weights $\pi_n$ are not the ordinary masses of the rendered nodes.
Instead, they are what we will call the \texttt{selection\_mass} values, because
they represent the user's probability of selecting those nodes.

Let the ordinary mass of visible node $n$ be
$$
m_n = \exp(z_n - z_r),
$$
where $r$ is the rendered root. Let $C(n)$ be the set of directly rendered
children of $n$, meaning those rendered descendants whose nearest rendered
ancestor is $n$. Then the selection\_mass of $n$ is
$$
\pi_n = m_n - \sum_{c \in C(n)} m_c.
$$
Equivalently, $\pi_n$ is the residual mass that stops at $n$ rather than
continuing into one of its directly rendered children. This is the correct
weight for timer placement because it is exactly the user's probability of
selecting node $n$.

\section{User Parameters}
\subsection{The Likelihood Model}
% We identify two random variables in the user's behavior: $T$, the "scan time" required for
% the user to determine their next signal, and $X$ the duration between intervals itself.
We model our observations of the user's signal as being subject to two sources of noise: an outlier probability, $\rho$,
that the signal was unintended, and additive noise on intended signals.
We model the unintended signals as uniformly distributed, and the noise as a Gaussian. (We do not require that the noise have
zero mean, since the user may be consistently late or early, relative to the intended time.) Hence our likelihood model for the
received signal is a uniform-gaussian mixture model, characterized by the outlier probability, 
$\rho$, and the parameters of the noise, $\mu_\epsilon$, $\sigma_\epsilon$.

\subsection{Directional Statistics}
% Should be a uniform-von_mises mixture model
% which we can estimate using the EM algorithm
% which consists in soft assignment / assignment of responsibilities


\subsection{A Hierarchical Temporal Model of User Statistics}
Performance will vary between users, between weeks for a given user,
between sessions for a given week, and between strings in a given session.

Here, I describe a simple, normal hierarchical model of user statistics.

\begin{figure}[H]
\centering
\makebox[\textwidth][c]{%
\begin{tikzpicture}[
  x=1.8cm,
  y=1.9cm,
  node/.style={
    draw,
    minimum width=2.5em,
    minimum height=1.5em,
    inner sep=4pt,
  },
  callout/.style={
    ->,
    very thick,
    dashed,
    draw=black!60,
  }
]
\node[node] (a1) at (-3.5, 3) {$c(0,0,0)$};
\node[node] (a2) at (-2.5, 3) {$c(0,0,1)$};
\node[node] (a3) at (-1.5, 3) {$c(0,1,0)$};
\node[node] (a4) at (-0.5, 3) {$c(0,1,1)$};
\node[node] (a5) at ( 0.5, 3) {$c(1,0,0)$};
\node[node] (a6) at ( 1.5, 3) {$c(1,0,1)$};
\node[node] (a7) at ( 2.5, 3) {$c(1,1,0)$};
\node[node] (a8) at ( 3.5, 3) {$c(1,1,1)$};

\node[node] (b1) at (-3, 2) {$c(i,0,0)$};
\node[node] (b2) at (-1, 2) {$c(i,0,1)$};
\node[node] (b3) at ( 1, 2) {$c(i,1,0)$};
\node[node] (b4) at ( 3, 2) {$c(i,1,1)$};

\node[node] (c1) at (-2, 1) {$c(i,j,0)$};
\node[node] (c2) at ( 2, 1) {$c(i,j,1)$};

\node[node] (d1) at (0, 0) {$X_{i,j,n}$};

\node[anchor=east] at (-4.4, 0) {Gesture};
\node[anchor=east] at (-4.4, 1) {String};
\node[anchor=east] at (-4.4, 2) {Session};
\node[anchor=east] at (-4.4, 3) {User};

\node[font=\normalsize] at ($(a2.center) + (0.18, -0.30)$) {exp};
\node[font=\normalsize] at ($(a4.center) + (0.18, -0.30)$) {exp};
\node[font=\normalsize] at ($(a6.center) + (0.18, -0.30)$) {exp};
\node[font=\normalsize] at ($(a8.center) + (0.18, -0.30)$) {exp};
\node[font=\normalsize] at ($(b2.center) + (0.18, -0.30)$) {exp};
\node[font=\normalsize] at ($(b4.center) + (0.18, -0.30)$) {exp};
\node[font=\normalsize] at ($(c2.center) + (0.18, -0.30)$) {exp};

\node[anchor=west, align=left, text width=3.4cm] (typical-sigma-label) at ($(c2.east) + (1.45, -0.18)$) {typical $\sigma_{\mathrm{string}}$\\for this session};
\draw[callout] (typical-sigma-label.west) to[out=180, in=0] (b3.east);

\node[anchor=west, align=left, text width=3.4cm] (session-control-label) at ($(c2.east) + (1.45, 0.55)$) {How much does this session control $\sigma_{\mathrm{string}}$?};
\draw[callout] (session-control-label.west) to[out=180, in=0] (b4.east);

\node[anchor=west, align=left, text width=3.4cm] (sigma-label) at ($(c2.east) + (1.45, -0.70)$) {$\sigma$ for this string};
\draw[callout] (sigma-label.west) to[out=180, in=0] (c2.east);

\coordinate (merge-b1) at ($(b1.north) + (0, 0.35)$);
\coordinate (merge-b2) at ($(b2.north) + (0, 0.35)$);
\coordinate (merge-b3) at ($(b3.north) + (0, 0.35)$);
\coordinate (merge-b4) at ($(b4.north) + (0, 0.35)$);
\coordinate (merge-c1) at ($(c1.north) + (0, 0.35)$);
\coordinate (merge-c2) at ($(c2.north) + (0, 0.35)$);
\coordinate (merge-d1) at ($(d1.north) + (0, 0.35)$);

\draw (a1.south) .. controls +(0, -0.25) and +(-0.22, 0.18) .. (merge-b1);
\draw (a2.south) .. controls +(0, -0.25) and +( 0.22, 0.18) .. (merge-b1);
\draw (merge-b1) -- node[pos=0.35, right] {$\mathcal N$} (b1.north);

\draw (a3.south) .. controls +(0, -0.25) and +(-0.22, 0.18) .. (merge-b2);
\draw (a4.south) .. controls +(0, -0.25) and +( 0.22, 0.18) .. (merge-b2);
\draw (merge-b2) -- node[pos=0.35, right] {$\mathcal N$} (b2.north);

\draw (a5.south) .. controls +(0, -0.25) and +(-0.22, 0.18) .. (merge-b3);
\draw (a6.south) .. controls +(0, -0.25) and +( 0.22, 0.18) .. (merge-b3);
\draw (merge-b3) -- node[pos=0.35, right] {$\mathcal N$} (b3.north);

\draw (a7.south) .. controls +(0, -0.25) and +(-0.22, 0.18) .. (merge-b4);
\draw (a8.south) .. controls +(0, -0.25) and +( 0.22, 0.18) .. (merge-b4);
\draw (merge-b4) -- node[pos=0.35, right] {$\mathcal N$} (b4.north);

\draw (b1.south) .. controls +(0, -0.25) and +(-0.22, 0.18) .. (merge-c1);
\draw (b2.south) .. controls +(0, -0.25) and +( 0.22, 0.18) .. (merge-c1);
\draw (merge-c1) -- node[pos=0.35, right] {$\mathcal N$} (c1.north);

\draw (b3.south) .. controls +(0, -0.25) and +(-0.22, 0.18) .. (merge-c2);
\draw (b4.south) .. controls +(0, -0.25) and +( 0.22, 0.18) .. (merge-c2);
\draw (merge-c2) -- node[pos=0.35, right] {$\mathcal N$} (c2.north);

\draw (c1.south) .. controls +(0, -0.25) and +(-0.22, 0.18) .. (merge-d1);
\draw (c2.south) .. controls +(0, -0.25) and +( 0.22, 0.18) .. (merge-d1);
\draw (merge-d1) -- node[pos=0.35, right] {$\mathcal N$} (d1.north);
\end{tikzpicture}
}
\caption{A hierarchical temporal model in which user-level statistics refine through session- and string-level variables to individual gestures.}
\label{fig:hierarchical-temporal-model}
\end{figure}

\begin{table}[h]
    \centering
    \renewcommand{\arraystretch}{1.4}
    \begin{tabular}{|c|c|c|}
        \hline
        \textbf{Level} & \textbf{Number of Latent Variables} & \textbf{Number of Data Points} \\
        \hline
        Gesture        & 0 & $N$ \\
        String         & $I \times J \times 2$ & 0 \\
        Session        & $I \times 4$ & 0 \\
        User           & 8 & 0 \\
        \hline
    \end{tabular}
    \caption{Number of latent variables and data points at each hierarchical level.}
\end{table}

In this model each data point is modelled as coming from 
a normal distribution parameterized by latent variables
specific to the narrowest temporal scope. Those variables
in turn are parameterized by variables of the next broader scope,
etc. In figure \ref{fig:hierarchical-temporal-model} 
a variable label that ends in $0$ indicates a mean for the 
next lower scope, while a variable label that ends in $1$ indicates
a standard deviation for the next lower scope.
Figure \ref{fig:hierarchical-temporal-model} also indicates
that deviation variables are in log-space, so the
distribution of a raw deviation variable is log-normal
with respect to fixed parent variables.

\subsubsection{Hierarchical Variational Inference}

We pursue estimation of these latent variables
using variational inference and a diagonal Gaussian
variational distribution. Evaluating the expected log likelihood
involves three closely related terms. We begin from
\[
C \mid A, B \sim \mathcal{N}(A, \exp(B)).
\]
It is convenient to evaluate this density in standardized coordinates,
\[
Z = \frac{C - A}{\exp(B)},
\]
which is simply a change of basis from the original variable $C$
to the corresponding standard-normal coordinate $Z$.
Under this change of basis,
\[
p(c \mid a, b)
= \frac{1}{\exp(b)} \mathcal{N}\left(
    \frac{c-a}{\exp(b)} ; 0, 1
\right),
\]
so in log space we must keep both the standardized Gaussian term
and the Jacobian term $-b$.
\begin{enumerate}
    \item \textbf{Latent Transition Likelihood}:
    \label{case:latent-transition}
    Given a latent variable $c$, and its two parent variables,
    $a$, and $b$, denote the variational distribution
    for these parameters as $(\mu_a, \sigma_a), (\mu_b, \sigma_b)$, and
    $(\mu_c, \sigma_c)$ respectively.
    Then our model that $C \sim \mathcal{N}(A, \exp(B))$ implies that
    the expected log likelihood for this transition is
    $$
    \mathbb{E}_{e_1,e_2,e_3} \left[
        \log \mathcal{N}\left(
            \frac{(\mu_c + e_3 \sigma_c) - (\mu_a + e_1 \sigma_a)}
            {\exp(\mu_b + e_2 \sigma_b)} ; 0, 1
        \right)
        - (\mu_b + e_2 \sigma_b)
    \right]
    $$
    where $e_1, e_2, e_3$ are iid standard normal random variables.
    \item \textbf{Initial Latent Likelihood}:
    In this case we consider that the parent variables are fixed and known
    (i.e. they are true prior parameters of the model, not latents)
    This is the same as case \ref{case:latent-transition}
    except that $a$ and $b$ are known constants. So we have
    an expected log likelihood of,
    $$
    \mathbb{E}_{e_3} \left[
        \log \mathcal{N}\left(
            \frac{(\mu_c + e_3 \sigma_c) - a}{\exp(b)} ; 0, 1
        \right)
        - b
    \right]
    $$
    \item \textbf{Data Likelihood}:
    This is the same as case \ref{case:latent-transition}
    except that we replace the child latent variable, $c$, with the known data point, $x$.
    So we have an expected log likelihood of,
    $$
    \mathbb{E}_{e_1,e_2} \left[
        \log \mathcal{N}\left(
            \frac{x - (\mu_a + e_1 \sigma_a)}{\exp(\mu_b + e_2 \sigma_b)} ; 0, 1
        \right)
        - (\mu_b + e_2 \sigma_b)
    \right]
    $$
\end{enumerate}
This suggests an implementation simplification: we need not make
any distinction between root prior parameters at the top level, intermediate latents,
and data points at the bottom level. 
We may calculate all the likelihood terms in the same way,
but enforce that in our variational distribution for a node $a$ with a known value,
whether a prior or a data point, we have that $\sigma_a = 0$
and $\mu_a$ is equal to the known value, with neither of the variables
being optimized.

Luckily, the above term may be evaluated analytically.
\begin{align}
    &\mathbb{E}_{e_1,e_2,e_3} \left[
        \log \mathcal{N}\left(
            \frac{(\mu_c + e_3 \sigma_c) - (\mu_a + e_1 \sigma_a)}
            {\exp(\mu_b + e_2 \sigma_b)} ; 0, 1
        \right)
        - (\mu_b + e_2 \sigma_b)
    \right] \\
    &= -\frac{1}{2}\log(2\pi)
    - \mathbb{E}_{e_2}[\mu_b + e_2 \sigma_b]
    - \frac{1}{2}
    \mathbb{E}_{e_1,e_2,e_3} \left[
        \left(
            \frac{(\mu_c + e_3 \sigma_c) - (\mu_a + e_1 \sigma_a)}
            {\exp(\mu_b + e_2 \sigma_b)}
        \right)^2
    \right] \\
    &= -\frac{1}{2}\log(2\pi)
    - \mu_b
    - \frac{1}{2}
    \mathbb{E}_{e_1,e_3} \left[
        \left(
            (\mu_c - \mu_a) + e_3 \sigma_c - e_1 \sigma_a
        \right)^2
    \right]
    \mathbb{E}_{e_2} \left[
        \exp(-2\mu_b - 2 e_2 \sigma_b)
    \right] \\
    &= -\frac{1}{2}\log(2\pi)
    - \mu_b
    - \frac{1}{2}
    \left(
        (\mu_c - \mu_a)^2 + \sigma_c^2 + \sigma_a^2
    \right)
    \exp(-2\mu_b + 2\sigma_b^2).
\end{align}


\subsection{Online Variational Estimation}
The above promises to quickly become a terrifying morass.
Here I discuss an online variational estimation algorithm
that is small and tractable.
We have the following desiderata,
\begin{enumerate}
    \item Use wrapped normal likelihoods
    \item Estimate the outlier rate, $\rho$
    \item Model variation between strings and between sessions
\end{enumerate}

We parameterize the user with three parameters:
$m$, $s$, and $\rho$.
We use a mean-field variational distribution for the parameters,
with six variational parameters:
\begin{align*}
    m &\sim \mathcal{N}(\mu_{m,q}, \sigma_{m,q}) \\
    \log s &\sim \mathcal{N}(\mu_{s,q}, \sigma_{s,q}) \\
    \rho &\sim Beta(\exp(a_q), \exp(b_q))
\end{align*}

Our likelihood is given by,
$$
P(x | m, s, \rho) = \rho \frac{1}{P} + (1 - \rho)\sum_{k \in \{-1, 0, 1\}} \frac{1}{\sqrt{2 \pi s}} \exp\left( -\frac{(x - m - k P)^2}{2 s} \right)
$$

Let $p$ denote the prior variational factors from the previous step and $q$ the updated factors.
Under mean field, write the (differential) entropy of $q$ as $H(q)$ and the cross-entropy of $q$ relative to $p$ as $H_x(q,p)$, so that $D_{\mathrm{KL}}(q\Vert p) = H_x(q,p) - H(q)$.
The online objective is
\begin{equation*}
    J(q) = \mathbb{E}_{q}\bigl[\log P(x \mid \theta)\bigr] + H(q) - \lambda\, H_x(q,p),
\end{equation*}
where $\lambda \in (0,1]$ is the \texttt{cross\_entropy\_discount} on the cross-entropy toward $p$.
When $\lambda = 1$, this is $\mathbb{E}_q[\log P(x\mid\theta)] - D_{\mathrm{KL}}(q\Vert p)$; when $\lambda < 1$, only the alignment term $H_x(q,p)$ is relaxed, while the entropy bonus $H(q)$ remains at full strength.

It is important to be explicit about what was wrong with an earlier shortcut: \emph{discounting the full $\mathrm{KL}(q\Vert p)$} incorrectly scales both the entropy and the cross-entropy pieces together.
That weakens proper entropy regularization in the same proportion as forgetting, which is not what we want.
\emph{Discounting only the cross-entropy term} $H_x(q,p)$ is the correct way to implement a forgetting heuristic while keeping full entropy regularization.

Because the factors are independent,
\begin{align*}
    H(q) &= H(q_m) + H(q_{\log s}) + H(q_\rho), \\
    H_x(q,p) &= H_x(q_m, p_m) + H_x(q_{\log s}, p_{\log s}) + H_x(q_\rho, p_\rho).
\end{align*}
For $m \sim \mathcal{N}(\mu_{m,q}, \sigma_{m,q})$ in the same standard-deviation parameterization as above (and identically for $\log s$ with $(\mu_{s,\cdot}, \sigma_{s,\cdot})$),
\begin{align*}
    H(q_m) &= \ln \sigma_{m,q} + \tfrac{1}{2}\ln(2\pi e), \\
    H_x(q_m, p_m) &= \ln \sigma_{m,p} + \tfrac{1}{2}\ln(2\pi) + \frac{\sigma_{m,q}^2 + (\mu_{m,q} - \mu_{m,p})^2}{2\sigma_{m,p}^2}.
\end{align*}
For $\rho \sim \mathrm{Beta}(\alpha_q, \beta_q)$ with $\alpha_q = \exp(a_q)$, $\beta_q = \exp(b_q)$ (and similarly $p$),
\begin{align*}
    H(q_\rho) &= \ln B(\alpha_q, \beta_q) - (\alpha_q - 1)\psi(\alpha_q) - (\beta_q - 1)\psi(\beta_q) + (\alpha_q + \beta_q - 2)\psi(\alpha_q + \beta_q), \\
    H_x(q_\rho, p_\rho) &= \ln B(\alpha_p, \beta_p) - (\alpha_p - 1)\mathbb{E}_q[\ln \rho] - (\beta_p - 1)\mathbb{E}_q[\ln(1-\rho)],
\end{align*}
where $B$ is the Beta function, $\psi$ is the digamma function, and $\mathbb{E}_q[\ln \rho] = \psi(\alpha_q) - \psi(\alpha_q + \beta_q)$, $\mathbb{E}_q[\ln(1-\rho)] = \psi(\beta_q) - \psi(\alpha_q + \beta_q)$.

We will approximate optimizing the log likelihood by optimizing
a lower bound derived by using Jensen's inequality
to pull the `logaddexp' out of the expectation.

\begin{align*}
    \mathbb{E}_{q} \left[ \log P(x | m, s, \rho) \right] &= \mathbb{E}_{q} \left[ \log \left( \rho \frac{1}{P} + (1 - \rho)\sum_{k \in \{-1, 0, 1\}} \frac{1}{\sqrt{2 \pi s}} \exp\left( -\frac{(x - m - k P)^2}{2 s} \right) \right) \right] \\
    &\geq \log \left( \exp\left( \mathbb{E}_{q} \left[ \log \rho \right] - \log P \right) + \sum_{k \in \{-1, 0, 1\}} \exp\left( \mathbb{E}_{q} \left[ \log(1 - \rho) - \frac{1}{2}\log(2\pi s) - \frac{(x - m - k P)^2}{2 s} \right] \right) \right)
\end{align*}

which is analytically tractable.
We will use the facts that
\begin{align*}
    \mathbb{E}_q \left[ \log \rho \right] &= \psi(\alpha_q) - \psi(\alpha_q + \beta_q) \\
    \mathbb{E}_q \left[ \log(1 - \rho) \right] &= \psi(\beta_q) - \psi(\alpha_q + \beta_q)
\end{align*}

Using the shorthand $\theta$ to denote the parameters $(m, \log s, \rho)$, combining the Jensen lower bound on $\mathbb{E}_q[\log P(x\mid\theta)]$ with the closed forms for $H(q)$ and $H_x(q,p)$ yields a tractable surrogate to maximize.

To make the expected likelihood component concrete, we substitute our earlier derivations into the lower bound.
We use the fact that if $\log s \sim \mathcal{N}(\mu_{s,q}, \sigma_{s,q})$, then $\mathbb{E}_q[s^{-1}] = \exp(-\mu_{s,q} + \sigma_{s,q}^2 / 2)$.
Thus, a fully broken-down lower bound on $J(q)$ is given by:
\begin{align*}
    J(q) &\geq \log \left( \exp\left( \psi(\alpha_q) - \psi(\alpha_q + \beta_q) - \log P \right) + \sum_{k \in \{-1, 0, 1\}} \exp\left( L_k \right) \right) \\
    &\quad + \ln \sigma_{m,q} + \tfrac{1}{2}\ln(2\pi e)
    + \ln \sigma_{s,q} + \tfrac{1}{2}\ln(2\pi e) \\
    &\quad + \ln B(\alpha_q, \beta_q) - (\alpha_q - 1)\psi(\alpha_q) - (\beta_q - 1)\psi(\beta_q) + (\alpha_q + \beta_q - 2)\psi(\alpha_q + \beta_q) \\
    &\quad - \lambda \Biggl(
        \ln \sigma_{m,p} + \tfrac{1}{2}\ln(2\pi) + \frac{\sigma_{m,q}^2 + (\mu_{m,q} - \mu_{m,p})^2}{2\sigma_{m,p}^2}
        + \ln \sigma_{s,p} + \tfrac{1}{2}\ln(2\pi) + \frac{\sigma_{s,q}^2 + (\mu_{s,q} - \mu_{s,p})^2}{2\sigma_{s,p}^2} \\
    &\qquad + \ln B(\alpha_p, \beta_p) - (\alpha_p - 1)\bigl(\psi(\alpha_q) - \psi(\alpha_q + \beta_q)\bigr)
        - (\beta_p - 1)\bigl(\psi(\beta_q) - \psi(\alpha_q + \beta_q)\bigr)
    \Biggr)
\end{align*}
where $L_k$ is the log-likelihood component for each wrapped normal mode:
\begin{align*}
    L_k = \psi(\beta_q) - \psi(\alpha_q + \beta_q) - \frac{1}{2}\log(2\pi) - \frac{1}{2}\mu_{s,q} - \frac{1}{2} \exp\left(-\mu_{s,q} + \frac{\sigma_{s,q}^2}{2}\right) \left((x - \mu_{m,q} - k P)^2 + \sigma_{m,q}^2\right)
\end{align*}

